% Introduction and Related Work draft for miniML
% This is a text fragment consumed by main.tex. It is NOT input via \input{}.

\section{Introduction}

Synaptic communication is the substrate for essentially every computation the brain performs, from sensory integration through working memory and motor control. Most of this communication is carried by discrete synaptic events: brief, stereotyped postsynaptic responses that follow the fusion of single neurotransmitter-filled vesicles \citep{kaeser2014release}. A sizeable fraction of these events occur spontaneously, independent of presynaptic action potentials, and these ``miniature'' events are far more than background noise. Their amplitude, frequency, and kinetics are a direct readout of receptor composition, presynaptic release probability, and the functional state of a synapse \citep{kavalali2015mechanisms, huganir2013ampars, reese2015spontaneous}. Changes in miniature-event statistics accompany homeostatic scaling, Hebbian plasticity, and disease states across model systems, and their quantitative analysis underpins much of modern cellular neurophysiology.

Despite their centrality, detecting and quantifying individual synaptic events in raw recordings remains stubbornly difficult. The events themselves are small relative to background noise, their timing is stochastic, and their waveform depends on the recording configuration, the receptor complement, the developmental stage of the preparation, and even the sampling hardware. The dominant algorithmic responses to these difficulties fall into three families. Finite-difference methods detect threshold crossings in the raw trace or in its derivative; template-matching methods convolve the trace with a stereotyped event waveform and look for peaks in the resulting score \citep{clements1997detection}; deconvolution methods invert a template to concentrate each event into a narrow impulse that is easier to threshold \citep{pernia2012deconvolution}. Bayesian inference procedures treat event times and amplitudes as latent variables and estimate them jointly under a generative model \citep{merel2016bayesian}. Each of these approaches works well in some regime and fails in another. Finite-threshold methods are simple but exquisitely sensitive to noise. Template-based methods require the user to supply (or tune) a waveform that closely matches the real events, and their output depends strongly on a user-chosen threshold. Bayesian methods are powerful but computationally expensive and require careful prior specification. In practice, many laboratories still rely on manual inspection or proprietary software \citep{mori2024simplyfire} to curate detected events, which introduces investigator-dependent variability that is incompatible with large-scale, reproducible analysis.

Deep learning has reshaped the way noisy, high-dimensional biological signals are processed \citep{lecun2015deep}. Convolutional neural networks excel at one-dimensional signal classification, and sequence models such as long short-term memory networks can capture the temporal structure of the surrounding context. In neuroscience specifically, CNN-based approaches now dominate spike inference from calcium imaging \citep{rupprecht2021cascade}, and transfer learning from large pretrained models is rapidly becoming the standard recipe for adapting a network trained on abundant data to a new domain with only limited labeled examples \citep{theodoris2023transfer, yosinski2014transferable}. These same ingredients, end-to-end learning, transfer, and automatic feature extraction, are a natural fit for spontaneous synaptic event detection, yet the community lacks an open, systematically benchmarked tool that delivers them in a form a wet-lab electrophysiologist can actually use.

Here we introduce miniML, a supervised deep learning framework for the automated, reproducible, and highly generalizable detection of synaptic events. miniML couples a compact CNN-LSTM classifier to a sliding-window inference procedure. It is trained on roughly 30{,}000 labeled segments from mossy fiber to cerebellar granule cell (MF-GC) mEPSC recordings, and then either applied directly to similar data or rapidly adapted to new preparations via transfer learning with as few as 400 labeled events. Against a panel of six established methods on synthetic ground-truth data, miniML delivers the highest F1 across the full range of signal-to-noise ratios encountered in real miniature event recordings, exhibits no false positives in the tested regime, and is virtually insensitive to threshold choice. It generalizes from mouse cerebellar slices to the Calyx of Held, human iPSC-derived cortical neurons, adult zebrafish telencephalon, Drosophila neuromuscular junctions, and optical imaging data acquired with both the glutamate sensor iGluSnFR3 \citep{aggarwal2023iglusnfr3} and the voltage indicator ASAP5 \citep{hao2024asap5}. In each of these settings, miniML either matches or exceeds the detection performance of the best alternative method, and in the low-SNR voltage-imaging regime it is the only method that recovers a substantial fraction of the ground-truth mEPSPs.

\section{Related Work}

\paragraph{Classical detection of spontaneous synaptic events.}
The two most widely used algorithms, sliding-template matching \citep{clements1997detection} and deconvolution \citep{pernia2012deconvolution}, have defined the field for more than two decades. Both require a fixed template that represents the expected event waveform. When the template is well matched to the data, they are sensitive and fast; when event kinetics drift because of a different preparation, receptor composition, or recording mode, their performance degrades and, critically, their output is strongly modulated by the chosen detection threshold. Finite-difference methods \citep{kudoh2002template} are simpler still but push the burden of noise rejection entirely onto the threshold. More recent tools such as SimplyFire \citep{mori2024simplyfire} provide usable GUIs and sensible defaults on top of these classical ideas, but they inherit the same structural limitations.

\paragraph{Probabilistic event inference.}
Bayesian approaches frame event detection as inference in a generative model of the recording \citep{merel2016bayesian}. This confers robustness to noise and a principled way to combine electrophysiological traces with auxiliary observations (for example, simultaneous optical reporters). The cost is computational: MCMC procedures are orders of magnitude slower than convolutional detectors, which makes them unattractive for analyzing the minutes-long recordings that are typical of miniature event experiments.

\paragraph{Deep learning for one-dimensional neurophysiological signals.}
Deep learning broke into neuroscience through imaging and spike inference \citep{lecun2015deep, rupprecht2021cascade}, and recent reviews of one-dimensional time-series classification show that CNN-LSTM architectures routinely outperform hand-engineered baselines on a wide variety of biological signals \citep{fawaz2019deeptsc}. Transfer learning, in which a model trained on a data-rich source task is partially reused on a target task, has become a standard recipe \citep{yosinski2014transferable, theodoris2023transfer}. miniML is, to our knowledge, the first open framework to bring this combination, an end-to-end CNN-LSTM classifier plus transfer learning, to spontaneous synaptic event detection, and to systematically benchmark it against both classical and probabilistic baselines on shared synthetic ground-truth data.

\paragraph{Optical detection of synaptic events.}
The recent generation of genetically encoded indicators, in particular iGluSnFR3 for glutamate \citep{aggarwal2023iglusnfr3} and ASAP5 for membrane voltage \citep{hao2024asap5}, has made it feasible to record unitary synaptic events with cellular and subcellular resolution. These recordings are, however, substantially noisier than patch-clamp electrophysiology and have a much lower sampling rate, which exposes the weaknesses of template-matching and deconvolution methods that were designed for high-SNR current traces. A detector that tolerates slower kinetics, lower SNR, and variable noise profiles is therefore an increasingly pressing need, and precisely the gap that miniML targets.
